\documentclass[12pt, a4paper]{article}
% ============================================================
%  preamble.tex — تنظیمات مشترک اسناد LaTeX (کامپایل با XeLaTeX)
% ============================================================
\usepackage{xcolor}
\usepackage{geometry}
\geometry{margin=2.5cm}
\usepackage{titlesec}
\usepackage{enumitem}
\setlist{nosep, leftmargin=*}
\usepackage{longtable}
\usepackage{booktabs}
\usepackage{array}
\usepackage{calc}
\usepackage{ragged2e}
\usepackage{fancyhdr}
\usepackage{etoolbox}
\usepackage{footnote}

% ---------- لازمهٔ جدول‌های خروجی pandoc ----------
\makeatletter
\patchcmd\longtable{\par}{\if@noskipsec\mbox{}\fi\par}{}{}
\makeatother
\makesavenoteenv{longtable}
\providecommand{\tightlist}{\setlength{\itemsep}{0pt}\setlength{\parskip}{0pt}}

% ---------- ظاهر ----------
\definecolor{primary}{RGB}{31, 78, 121}
\definecolor{accent}{RGB}{0, 150, 136}
\titleformat{\section}{\Large\bfseries\color{primary}}{\thesection}{1em}{}
\titleformat{\subsection}{\large\bfseries}{\thesubsection}{1em}{}
\titleformat{\subsubsection}{\normalsize\bfseries}{\thesubsubsection}{1em}{}
\renewcommand{\arraystretch}{1.3}
\linespread{1.3}

\pagestyle{fancy}
\fancyhf{}
\fancyfoot[C]{\thepage}
\renewcommand{\headrulewidth}{0.4pt}
\setlength{\headheight}{14pt}

% ---------- hyperref قبل از xepersian ----------
\usepackage{hyperref}
\hypersetup{colorlinks=true, linkcolor=primary, urlcolor=primary, citecolor=primary}

% ---------- فارسی‌سازی (آخرین بسته) ----------
\usepackage{xepersian}
\settextfont{Tahoma}
\setlatintextfont{Times New Roman}
\setdigitfont{Tahoma}

% ---------- شمارندهٔ لازم برای جدول‌های pandoc (LTcaptype = none) ----------
\newcounter{none}

\fancyhead[R]{\small سند چشم‌انداز — کلاس‌یار}
\fancyhead[L]{\small نسخهٔ ۱.۰}

\begin{document}

\begin{titlepage}
\centering
\vspace*{3cm}
{\Huge\bfseries سند چشم‌انداز\par}
\vspace{0.5cm}
{\LARGE\bfseries پروژهٔ کلاس‌یار\par}
\vspace{0.5cm}
{\large پلتفرم اشتراک جزوه، نمونه‌سؤال و ارزیابی استادان\par}
\vspace{2cm}
{\large تهیه‌کنندگان:\par}
{\large مریم قندی زنجانی، ملیکا شهبازی، آنیسا فروغمند، مائده رضایی، محمدمهدی ذوالقدر\par}
\vspace{0.5cm}
{\large دانشگاه آزاد زنجان\par}
\vspace{1cm}
{\large ۱۴۰۵/۰۷/۱۱\par}
\end{titlepage}

\tableofcontents
\newpage

{\def\LTcaptype{none} % do not increment counter
\begin{longtable}[]{@{}
  >{\raggedright\arraybackslash}p{(\linewidth - 2\tabcolsep) * \real{0.5000}}
  >{\raggedright\arraybackslash}p{(\linewidth - 2\tabcolsep) * \real{0.5000}}@{}}
\toprule\noalign{}
\endhead
\bottomrule\noalign{}
\endlastfoot
\textbf{نام پروژه} & کلاس‌یار (ClassYar) — پلتفرم اشتراک جزوه، نمونه‌سؤال و ارزیابی استادان \\
\textbf{نسخهٔ سند} & 1.0 \\
\textbf{تاریخ} & ۱۴۰۵/۰۷/۱۱ \\
\textbf{تهیه‌کنندگان} & {[}مریم قندی زنجانی، ملیکا شهبازی،آنیسا فروغمند، مائده رضایی، محمدمهدی ذوالقدر{]} \\
\textbf{وضعیت} & پیش‌نویس \\
\textbf{آخرین به‌روزرسانی} & ۱۴۰۵/۰۷/۱۱ \\
\end{longtable}
}

\begin{center}\rule{0.5\linewidth}{0.5pt}\end{center}

\section{مقدمه}

\subsection{هدف سند}

این سند چشم‌انداز، دامنه و ذی‌نفعان پروژهٔ \textbf{کلاس‌یار} را تعریف می‌کند.
هدف آن ایجاد یک درک مشترک میان تیم توسعه، کاربران و ذی‌نفعان درباره
«چه چیزی می‌سازیم، برای چه کسی، و چرا» است.

این سند \textbf{قرارداد رسمی} نیست، بلکه \textbf{مرجع راهنما} برای تصمیم‌گیری‌های
بعدی در طول چرخهٔ حیات پروژه است و در هر Sprint بازبینی می‌شود.

\subsection{دامنهٔ سند}

\begin{itemize}
\tightlist
\item
  \textbf{در دامنه:} اهداف محصول، کاربران هدف، محدودهٔ MVP، معیارهای موفقیت،
  ریسک‌های سطح بالا.
\item
  \textbf{خارج از دامنه:} جزئیات فنی، طراحی UI، نیازمندی‌های دقیق functional (که
  در SRS می‌آید).
\end{itemize}

\subsection{تعاریف و اصطلاحات}

{\def\LTcaptype{none} % do not increment counter
\begin{longtable}[]{@{}
  >{\raggedright\arraybackslash}p{(\linewidth - 2\tabcolsep) * \real{0.5000}}
  >{\raggedright\arraybackslash}p{(\linewidth - 2\tabcolsep) * \real{0.5000}}@{}}
\toprule\noalign{}
\begin{minipage}[b]{\linewidth}\raggedright
اصطلاح
\end{minipage} & \begin{minipage}[b]{\linewidth}\raggedright
تعریف
\end{minipage} \\
\midrule\noalign{}
\endhead
\bottomrule\noalign{}
\endlastfoot
\textbf{MVP} & Minimum Viable Product — کمینهٔ محصول قابل ارائه \\
\textbf{User Story} & توصیف نیازمندی از دید کاربر \\
\textbf{Persona} & نمایندهٔ فرضی کاربر هدف \\
\textbf{Sprint} & بازهٔ کوتاه توسعه (۲-۳ هفته) در Scrum \\
\textbf{DoD} & Definition of Done — معیار «تمام‌شده» بودن کار \\
\textbf{جزوه} & هر فایل آموزشی دانشجویی مانند جزوهٔ دست‌نویس اسکن‌شده، تایپ‌شده یا خلاصهٔ درس \\
\textbf{واحد دانشگاهی} & شعبه یا دانشکدهٔ مشخص یک دانشگاه (مثلاً دانشگاه آزاد، واحد زنجان) \\
\textbf{نمرهٔ تشویقی} & امتیاز یا نمرهٔ اضافه‌ای که استاد برای مشارکت مفید دانشجو ثبت می‌کند \\
\textbf{MoSCoW} & روش اولویت‌بندی: Must، Should، Could، Won\textquotesingle t \\
\end{longtable}
}

\subsection{مراجع}

\begin{itemize}
\tightlist
\item
  IEEE 830 — Recommended Practice for Software Requirements Specifications
\item
  Roman Pichler — \emph{Product Vision Board}
\item
  Scrum Guide (2020)
\item
  Sommerville — \emph{Software Engineering}, Chapter 1-2
\end{itemize}

\begin{center}\rule{0.5\linewidth}{0.5pt}\end{center}

\section{چشم‌انداز محصول}

\subsection{بیانیهٔ چشم‌انداز (Vision Statement)}

\begin{quote}
\textbf{برای} دانشجویانی \textbf{که} برای پیدا کردن جزوه، نمونه‌سؤال و اطلاعات مطمئن
دربارهٔ استادان خود بین گروه‌های پیام‌رسان و فایل‌های پراکنده سرگردان‌اند،
\textbf{کلاس‌یار} یک \textbf{پلتفرم وب دانشجویی} است
\textbf{که} اشتراک جزوه، جست‌وجو بر پایهٔ استاد، استان و واحد دانشگاهی،
امتیازدهی و گفت‌وگو را در یک جا فراهم می‌کند.

\textbf{برخلاف} گروه‌های تلگرام و ایتا که محتوا در آن‌ها گم می‌شود و کیفیتش
مشخص نیست، \textbf{محصول ما} محتوا را بر پایهٔ استاد و درس سازمان می‌دهد و با
امتیاز دانشجویان و تأیید استادان قابل‌اعتماد می‌کند.
\end{quote}

\subsection{ارزش پیشنهادی (Value Proposition)}

{\def\LTcaptype{none} % do not increment counter
\begin{longtable}[]{@{}
  >{\raggedright\arraybackslash}p{(\linewidth - 2\tabcolsep) * \real{0.4615}}
  >{\raggedright\arraybackslash}p{(\linewidth - 2\tabcolsep) * \real{0.5385}}@{}}
\toprule\noalign{}
\begin{minipage}[b]{\linewidth}\raggedright
برای
\end{minipage} & \begin{minipage}[b]{\linewidth}\raggedright
ارزش
\end{minipage} \\
\midrule\noalign{}
\endhead
\bottomrule\noalign{}
\endlastfoot
\textbf{دانشجوی استفاده‌کننده} & پیدا کردن سریع جزوه و نمونه‌سؤال استاد موردنظر، با دیدن امتیاز و نظر دیگران قبل از دانلود \\
\textbf{دانشجوی مشارکت‌کننده} & اشتراک جزوه، دریافت بازخورد و امکان گرفتن نمرهٔ تشویقی از استاد \\
\textbf{استاد} & دیدن منابع دانشجویی مربوط به درس خود، تأیید جزوه‌های خوب و تشویق مشارکت‌کنندگان \\
\textbf{مدیر سامانه} & کنترل محتوا و کاربران و حفظ اعتبار پلتفرم \\
\end{longtable}
}

\subsection{مزیت رقابتی}

\begin{itemize}
\tightlist
\item
  تمرکز: همه‌چیز بر پایهٔ «استاد + درس + واحد دانشگاهی» سازمان می‌یابد، نه یک انبار فایل
\item
  اعتماد: امتیاز دانشجویان و تأیید استاد، کیفیت جزوه را نشان می‌دهد
\item
  انگیزه: استاد می‌تواند برای مشارکت مفید نمرهٔ تشویقی در نظر بگیرد
\item
  ارتباط: نظرها، لینک گروه‌های کلاسی و نمونه‌سؤال‌ها در کنار هم هستند
\item
  سادگی: فارسی، راست‌به‌چپ و بدون پیچیدگی
\end{itemize}

\begin{center}\rule{0.5\linewidth}{0.5pt}\end{center}

\section{مشکل و فرصت}

\subsection{مشکل فعلی (Problem Statement)}

دانشجویان و استادان با چالش‌های زیر روبه‌رو هستند:
1. \textbf{پراکندگی منابع:} جزوه و نمونه‌سؤال در گروه‌های پیام‌رسان و حافظهٔ شخصی افراد پخش است.
2. \textbf{سخت‌بودن پیدا کردن:} نمی‌توان جزوهٔ یک استاد مشخص در یک واحد دانشگاهی مشخص را فیلتر کرد.
3. \textbf{نبود معیار کیفیت:} دانشجو نمی‌داند کدام جزوه کامل و قابل‌اعتماد است.
4. \textbf{نبود اطلاعات دربارهٔ استاد:} نظر و تجربهٔ دانشجویان قبلی در دسترس نیست.
5. \textbf{نبود انگیزه برای اشتراک:} دانشجویی که جزوهٔ کامل دارد دلیلی برای به‌اشتراک‌گذاشتن آن ندارد.
6. \textbf{نبود ارتباط استاد با منابع دانشجویی:} استاد نمی‌داند دانشجویان از چه منابعی استفاده می‌کنند.

\subsection{فرصت (Opportunity)}

\begin{itemize}
\tightlist
\item
  نیاز مشترک میان دانشجویان همهٔ دانشگاه‌ها و رشته‌ها
\item
  امکان ساخت یک مخزن مشترک و قابل‌جست‌وجو از جزوه‌های واقعی
\item
  امکان ایجاد انگیزه با نمرهٔ تشویقی و امتیازدهی
\item
  ضعف گروه‌های پیام‌رسان در سازمان‌دهی محتوا
\item
  امکان توسعه به همهٔ استان‌ها و دانشگاه‌های کشور
\end{itemize}

\begin{center}\rule{0.5\linewidth}{0.5pt}\end{center}

\section{کاربران هدف و ذی‌نفعان}

\subsection{ذی‌نفعان (Stakeholders)}

{\def\LTcaptype{none} % do not increment counter
\begin{longtable}[]{@{}
  >{\raggedright\arraybackslash}p{(\linewidth - 4\tabcolsep) * \real{0.3810}}
  >{\raggedright\arraybackslash}p{(\linewidth - 4\tabcolsep) * \real{0.2857}}
  >{\raggedright\arraybackslash}p{(\linewidth - 4\tabcolsep) * \real{0.3333}}@{}}
\toprule\noalign{}
\begin{minipage}[b]{\linewidth}\raggedright
ذی‌نفع
\end{minipage} & \begin{minipage}[b]{\linewidth}\raggedright
نقش
\end{minipage} & \begin{minipage}[b]{\linewidth}\raggedright
علاقه
\end{minipage} \\
\midrule\noalign{}
\endhead
\bottomrule\noalign{}
\endlastfoot
دانشجوی استفاده‌کننده & جست‌وجو، دانلود، امتیاز و نظر & پیدا کردن سریع منبع باکیفیت \\
دانشجوی مشارکت‌کننده & آپلود جزوه و نمونه‌سؤال & دیده‌شدن، بازخورد و نمرهٔ تشویقی \\
استاد & تأیید جزوه، نمرهٔ تشویقی، مشاهدهٔ نظرات & منابع درست برای درس خود و تشویق دانشجویان \\
مدیر سامانه & مدیریت کاربران، اطلاعات پایه و گزارش‌ها & نظم، اعتبار و امنیت پلتفرم \\
تیم توسعه & سازندهٔ محصول & تحویل به‌موقع، یادگیری \\
استاد درس & ناظر پروژه & رعایت اصول مهندسی \\
\end{longtable}
}

\subsection{پرسوناهای کلیدی}

جزئیات کامل در \texttt{docs/02-requirements/personas.md}. خلاصه:

{\def\LTcaptype{none} % do not increment counter
\begin{longtable}[]{@{}
  >{\raggedright\arraybackslash}p{(\linewidth - 4\tabcolsep) * \real{0.2857}}
  >{\raggedright\arraybackslash}p{(\linewidth - 4\tabcolsep) * \real{0.2143}}
  >{\raggedright\arraybackslash}p{(\linewidth - 4\tabcolsep) * \real{0.5000}}@{}}
\toprule\noalign{}
\begin{minipage}[b]{\linewidth}\raggedright
پرسونا
\end{minipage} & \begin{minipage}[b]{\linewidth}\raggedright
نقش
\end{minipage} & \begin{minipage}[b]{\linewidth}\raggedright
ویژگی کلیدی
\end{minipage} \\
\midrule\noalign{}
\endhead
\bottomrule\noalign{}
\endlastfoot
\textbf{نگار} & دانشجوی استفاده‌کننده & ۲۰ ساله، وقت کم و نیازمند جزوهٔ مطمئن \\
\textbf{امیرحسین} & دانشجوی مشارکت‌کننده & ۲۲ ساله، جزوه‌های منظم دارد و دنبال انگیزه است \\
\textbf{دکتر کاظمی} & استاد & ۴۶ ساله، مهارت فنی متوسط، دنبال تشویق دانشجویان فعال \\
\textbf{خانم نوروزی} & مدیر سامانه & ۲۶ ساله، مسئول انجمن علمی، نمایندهٔ مشتری \\
\end{longtable}
}

\begin{center}\rule{0.5\linewidth}{0.5pt}\end{center}

\section{محدودهٔ محصول}

\subsection{داخل محدوده (In Scope — MVP)}

\textbf{کاربری (Functional):}
- \textbf{F-01} ثبت‌نام و ورود با نقش دانشجو یا استاد
- \textbf{F-02} مدیریت اطلاعات پایه (استان، دانشگاه، واحد، درس، استاد) توسط مدیر؛ دانشجو می‌تواند افزودن استاد را درخواست دهد
- \textbf{F-03} آپلود جزوه با مشخصات (استاد، درس، استان، واحد، ترم، توضیح)
- \textbf{F-04} جست‌وجو و فیلتر بر پایهٔ استاد، استان، واحد دانشگاهی، درس و ترم
- \textbf{F-05} صفحهٔ جزئیات جزوه و دانلود فایل
- \textbf{F-06} امتیازدهی ۱ تا ۵ به جزوه و نمایش میانگین
- \textbf{F-07} صفحهٔ استاد: جزوه‌ها، نمونه‌سؤال‌ها، لینک گروه‌ها و میانگین امتیاز
- \textbf{F-08} امتیازدهی دانشجو به استاد
- \textbf{F-09} نظرها روی جزوه و صفحهٔ استاد، همراه با امکان پاسخ
- \textbf{F-10} ثبت و نمایش لینک گروه‌های کلاسی و نمونه‌سؤال‌ها
- \textbf{F-11} تأیید جزوه و ثبت نمرهٔ تشویقی توسط استاد
- \textbf{F-12} گزارش محتوای نامناسب یا ناقض حقوق دیگران
- \textbf{F-13} پنل مدیر: کاربران، گزارش‌ها و حذف محتوا

\textbf{غیرکاربری (Non-functional):}
- \textbf{NF-01} زمان بارگذاری صفحه \textless{} ۲ ثانیه
- \textbf{NF-02} پشتیبانی از ۱۰۰ کاربر همزمان
- \textbf{NF-03} رابط کاربری واکنش‌گرا (Responsive) برای موبایل و دسکتاپ
- \textbf{NF-04} رمزنگاری گذرواژه‌ها (bcrypt)
- \textbf{NF-05} پشتیبانی از زبان فارسی و RTL
- \textbf{NF-06} اعتبارسنجی فایل‌ها (فقط PDF، JPG، PNG و DOCX، حداکثر ۲۰ مگابایت)
- \textbf{NF-07} کنترل دسترسی بر پایهٔ نقش (دانشجو، استاد، مدیر)
- \textbf{NF-08} هر کاربر برای هر جزوه و هر استاد فقط یک امتیاز بتواند ثبت کند

\subsection{خارج از محدوده (Out of Scope)}

{\def\LTcaptype{none} % do not increment counter
\begin{longtable}[]{@{}
  >{\raggedright\arraybackslash}p{(\linewidth - 2\tabcolsep) * \real{0.4615}}
  >{\raggedright\arraybackslash}p{(\linewidth - 2\tabcolsep) * \real{0.5385}}@{}}
\toprule\noalign{}
\begin{minipage}[b]{\linewidth}\raggedright
مورد
\end{minipage} & \begin{minipage}[b]{\linewidth}\raggedright
دلیل
\end{minipage} \\
\midrule\noalign{}
\endhead
\bottomrule\noalign{}
\endlastfoot
ثبت نمرهٔ رسمی در سامانهٔ دانشگاه & نمرهٔ تشویقی فقط در کلاس‌یار ثبت می‌شود و اتصال به سامانهٔ آموزشی خارج از توان یک ترم است \\
اپلیکیشن موبایل & تمرکز روی وب \\
چت زنده بین کاربران & نظرها و پاسخ‌ها کافی است \\
فروش جزوه و پرداخت & خارج از هدف پروژه \\
تشخیص خودکار سرقت یا کپی‌رایت & پیچیدگی فنی؛ فقط گزارش دستی \\
تبدیل جزوهٔ دست‌نویس به متن (OCR) & خارج از توان یک ترم \\
احراز هویت با سامانهٔ دانشگاه & نیاز به مجوز؛ فقط ایمیل و تأیید دستی استاد \\
چند زبانه (فارسی/انگلیسی) & تمرکز روی فارسی \\
\end{longtable}
}

\subsection{محدودیت‌ها (Constraints)}

\begin{itemize}
\tightlist
\item
  \textbf{زمانی:} ۱۴ هفته، ۳ Sprint
\item
  \textbf{تیمی:} ۳-۵ نفر
\item
  \textbf{فنی:} بدون سرور اختصاصی، استقرار روی سرویس ابری رایگان
\item
  \textbf{بودجه:} صفر (فقط ابزارهای رایگان)
\item
  \textbf{محتوایی:} مسئولیت محتوای بارگذاری‌شده با کاربر است و باید قوانین استفاده تعریف شود
\end{itemize}

\begin{center}\rule{0.5\linewidth}{0.5pt}\end{center}

\section{معیارهای موفقیت (Success Metrics)}

\subsection{معیارهای کاربری}

{\def\LTcaptype{none} % do not increment counter
\begin{longtable}[]{@{}
  >{\raggedright\arraybackslash}p{(\linewidth - 4\tabcolsep) * \real{0.2059}}
  >{\raggedright\arraybackslash}p{(\linewidth - 4\tabcolsep) * \real{0.1765}}
  >{\raggedright\arraybackslash}p{(\linewidth - 4\tabcolsep) * \real{0.6176}}@{}}
\toprule\noalign{}
\begin{minipage}[b]{\linewidth}\raggedright
معیار
\end{minipage} & \begin{minipage}[b]{\linewidth}\raggedright
هدف
\end{minipage} & \begin{minipage}[b]{\linewidth}\raggedright
نحوهٔ اندازه‌گیری
\end{minipage} \\
\midrule\noalign{}
\endhead
\bottomrule\noalign{}
\endlastfoot
زمان پیدا کردن جزوهٔ یک استاد مشخص & \textless{} ۲ دقیقه & تست کاربری \\
زمان آپلود یک جزوه & \textless{} ۳ دقیقه & تست کاربری \\
موفقیت در ثبت امتیاز و نظر & \textgreater{} ۹۰٪ کاربران آزمایشی & تست کاربری \\
رضایت کاربر & \textgreater{} ۸۰٪ & نظرسنجی ساده \\
\end{longtable}
}

\subsection{معیارهای فنی}

{\def\LTcaptype{none} % do not increment counter
\begin{longtable}[]{@{}ll@{}}
\toprule\noalign{}
معیار & هدف \\
\midrule\noalign{}
\endhead
\bottomrule\noalign{}
\endlastfoot
پوشش تست (Code Coverage) & \textgreater{} ۷۰٪ \\
تعداد باگ بحرانی در انتشار & صفر \\
زمان Build در CI & \textless{} ۵ دقیقه \\
موفقیت Deploy & ۱۰۰٪ \\
صحت کنترل دسترسی نقش‌ها & ۱۰۰٪ سناریوهای اصلی \\
\end{longtable}
}

\subsection{معیارهای فرآیندی}

\begin{itemize}
\tightlist
\item
  تحویل هر Sprint به‌موقع
\item
  رعایت Definition of Done در هر Story
\item
  مستندات به‌روز در پایان هر Sprint
\end{itemize}

\begin{center}\rule{0.5\linewidth}{0.5pt}\end{center}

\section{ریسک‌های سطح بالا}

{\def\LTcaptype{none} % do not increment counter
\begin{longtable}[]{@{}
  >{\raggedright\arraybackslash}p{(\linewidth - 6\tabcolsep) * \real{0.1818}}
  >{\raggedright\arraybackslash}p{(\linewidth - 6\tabcolsep) * \real{0.2424}}
  >{\raggedright\arraybackslash}p{(\linewidth - 6\tabcolsep) * \real{0.1515}}
  >{\raggedright\arraybackslash}p{(\linewidth - 6\tabcolsep) * \real{0.4242}}@{}}
\toprule\noalign{}
\begin{minipage}[b]{\linewidth}\raggedright
ریسک
\end{minipage} & \begin{minipage}[b]{\linewidth}\raggedright
احتمال
\end{minipage} & \begin{minipage}[b]{\linewidth}\raggedright
اثر
\end{minipage} & \begin{minipage}[b]{\linewidth}\raggedright
راهکار کاهش
\end{minipage} \\
\midrule\noalign{}
\endhead
\bottomrule\noalign{}
\endlastfoot
تغییر نیازمندی‌ها در میانهٔ پروژه & بالا & متوسط & Sprintهای کوتاه، Backlog منعطف \\
کمبود زمان به‌دلیل وسعت پروژه & بالا & بالا & MVP کوچک، اولویت MoSCoW؛ F-10 و F-11 اول حذف می‌شوند \\
Scope Creep & بالا & بالا & مرزبندی دقیق MVP، تأیید استاد \\
نظرهای توهین‌آمیز یا ناعادلانه دربارهٔ استادان & متوسط & بالا & گزارش محتوا، امکان پاسخ استاد، قوانین استفاده، حذف توسط مدیر \\
آپلود محتوای دارای حق مالکیت دیگران & متوسط & بالا & تعهد آپلودکننده، گزارش و حذف سریع، قوانین سامانه \\
امتیازدهی تقلبی با حساب‌های متعدد & متوسط & متوسط & یک امتیاز برای هر کاربر، تأیید ایمیل \\
مشارکت‌نکردن استادان در ابتدا & بالا & متوسط & سامانه بدون مشارکت استاد هم کار می‌کند؛ استاد فقط ارزش اضافه می‌دهد \\
ضعف فنی تیم در بخشی & متوسط & متوسط & منتورینگ استاد، تقسیم کار \\
از دست دادن داده & کم & بالا & Backup خودکار، Git \\
\end{longtable}
}

\begin{center}\rule{0.5\linewidth}{0.5pt}\end{center}

\section{رویکرد و متدولوژی}

\subsection{متدولوژی}

\begin{itemize}
\tightlist
\item
  \textbf{Scrum سبک‌شده:} ۳ Sprint دو-سه هفته‌ای
\item
  \textbf{ابزار مدیریت:} GitHub Projects
\item
  \textbf{جلسات:} Sprint Planning، Daily (کوتاه، آنلاین)، Review، Retrospective
\end{itemize}

\subsection{فناوری پیشنهادی}

{\def\LTcaptype{none} % do not increment counter
\begin{longtable}[]{@{}
  >{\raggedright\arraybackslash}p{(\linewidth - 2\tabcolsep) * \real{0.4286}}
  >{\raggedright\arraybackslash}p{(\linewidth - 2\tabcolsep) * \real{0.5714}}@{}}
\toprule\noalign{}
\begin{minipage}[b]{\linewidth}\raggedright
لایه
\end{minipage} & \begin{minipage}[b]{\linewidth}\raggedright
فناوری
\end{minipage} \\
\midrule\noalign{}
\endhead
\bottomrule\noalign{}
\endlastfoot
Backend & Node.js + Express \\
Frontend & React \\
Database & PostgreSQL \\
ORM & Prisma \\
ذخیرهٔ فایل & دیسک سرور یا سرویس ذخیرهٔ ابری رایگان \\
Version Control & Git + GitHub \\
CI/CD & GitHub Actions \\
مستندسازی & Markdown در GitHub \\
\end{longtable}
}

\begin{quote}
دلیل انتخاب: یک زبان (JavaScript) برای کل تیم، ابزارهای رایگان و مستندات فراوان،
و پایگاه دادهٔ رابطه‌ای برای ارتباط‌های استاد، درس، واحد و جزوه.
\end{quote}

\subsection{Definition of Done}

یک Story وقتی «تمام» است که:
1. کد نوشته و Merge شده باشد.
2. تست Unit برای آن نوشته شده باشد.
3. Code Review شده باشد.
4. مستندات مربوطه به‌روز شده باشد.
5. در محیط Demo کار کند.

\begin{center}\rule{0.5\linewidth}{0.5pt}\end{center}

\section{نقشهٔ راه (Roadmap)}

{\def\LTcaptype{none} % do not increment counter
\begin{longtable}[]{@{}
  >{\raggedright\arraybackslash}p{(\linewidth - 6\tabcolsep) * \real{0.2759}}
  >{\raggedright\arraybackslash}p{(\linewidth - 6\tabcolsep) * \real{0.2069}}
  >{\raggedright\arraybackslash}p{(\linewidth - 6\tabcolsep) * \real{0.2414}}
  >{\raggedright\arraybackslash}p{(\linewidth - 6\tabcolsep) * \real{0.2759}}@{}}
\toprule\noalign{}
\begin{minipage}[b]{\linewidth}\raggedright
Sprint
\end{minipage} & \begin{minipage}[b]{\linewidth}\raggedright
بازه
\end{minipage} & \begin{minipage}[b]{\linewidth}\raggedright
تمرکز
\end{minipage} & \begin{minipage}[b]{\linewidth}\raggedright
خروجی
\end{minipage} \\
\midrule\noalign{}
\endhead
\bottomrule\noalign{}
\endlastfoot
Sprint 1 & هفته ۱-۵ & پایه و کشف & ثبت‌نام و ورود، اطلاعات پایه، آپلود جزوه، فهرست و جزئیات جزوه (F-01 تا F-05) \\
Sprint 2 & هفته ۶-۱۱ & ارزیابی و تعامل & فیلتر کامل، امتیاز جزوه، صفحهٔ استاد، امتیاز استاد، نظرها (F-04 تا F-09) \\
Sprint 3 & هفته ۱۲-۱۴ & استاد، کیفیت و انتشار & لینک گروه و نمونه‌سؤال، تأیید و نمرهٔ تشویقی، گزارش و پنل مدیر، تست، CI/CD، Demo Day (F-10 تا F-13) \\
\end{longtable}
}

\begin{center}\rule{0.5\linewidth}{0.5pt}\end{center}

\section{تأیید و تصویب}

{\def\LTcaptype{none} % do not increment counter
\begin{longtable}[]{@{}
  >{\raggedright\arraybackslash}p{(\linewidth - 6\tabcolsep) * \real{0.2308}}
  >{\raggedright\arraybackslash}p{(\linewidth - 6\tabcolsep) * \real{0.2308}}
  >{\raggedright\arraybackslash}p{(\linewidth - 6\tabcolsep) * \real{0.2692}}
  >{\raggedright\arraybackslash}p{(\linewidth - 6\tabcolsep) * \real{0.2692}}@{}}
\toprule\noalign{}
\begin{minipage}[b]{\linewidth}\raggedright
نقش
\end{minipage} & \begin{minipage}[b]{\linewidth}\raggedright
نام
\end{minipage} & \begin{minipage}[b]{\linewidth}\raggedright
تاریخ
\end{minipage} & \begin{minipage}[b]{\linewidth}\raggedright
امضا
\end{minipage} \\
\midrule\noalign{}
\endhead
\bottomrule\noalign{}
\endlastfoot
صاحب محصول & دکترابراهیم فضلی (استاد درس) & {[}۱۴۰۵/۰۷/۱۱{]} & \\
مشتری فرضی & خانم نوروزی & {[}۱۴۰۵/۰۷/۱۱{]} & \\
نمایندهٔ تیم & {[}مریم قندی زنجانی{]} & {[}تاریخ{]} & \\
\end{longtable}
}

\begin{center}\rule{0.5\linewidth}{0.5pt}\end{center}

\section*{پیوست الف — تاریخچهٔ تغییرات}\addcontentsline{toc}{section}{پیوست الف — تاریخچهٔ تغییرات}

{\def\LTcaptype{none} % do not increment counter
\begin{longtable}[]{@{}
  >{\raggedright\arraybackslash}p{(\linewidth - 6\tabcolsep) * \real{0.1842}}
  >{\raggedright\arraybackslash}p{(\linewidth - 6\tabcolsep) * \real{0.1842}}
  >{\raggedright\arraybackslash}p{(\linewidth - 6\tabcolsep) * \real{0.2632}}
  >{\raggedright\arraybackslash}p{(\linewidth - 6\tabcolsep) * \real{0.3684}}@{}}
\toprule\noalign{}
\begin{minipage}[b]{\linewidth}\raggedright
نسخه
\end{minipage} & \begin{minipage}[b]{\linewidth}\raggedright
تاریخ
\end{minipage} & \begin{minipage}[b]{\linewidth}\raggedright
تغییرات
\end{minipage} & \begin{minipage}[b]{\linewidth}\raggedright
تهیه‌کنندگان
\end{minipage} \\
\midrule\noalign{}
\endhead
\bottomrule\noalign{}
\endlastfoot
1.0 & ۱۴۰۵/۰۷/۱۱ & نسخهٔ اولیه & {[}مریم قندی زنجانی، ملیکا شهبازی،آنیسا فروغمند، مائده رضایی، محمدمهدی ذوالقدر{]} \\
\end{longtable}
}

\end{document}
